\section{USAJMO 2017}
        \begin{enumerate}
        	\item (\textit{USAJMO 2017}) Note: For any geometry problem whose statement begins with an asterisk ($*$), the first page of the solution must be a large, in-scale, clearly labeled diagram. Failure to meet this requirement will result in an automatic 1-point deduction. 


 Prove that there are infinitely many distinct pairs $(a,b)$ of relatively prime positive integers $a > 1$ and $b > 1$ such that $a^b + b^a$ is divisible by $a + b.$


 

	\item (\textit{USAJMO 2017}) Note: For any geometry problem whose statement begins with an asterisk ($*$), the first page of the solution must be a large, in-scale, clearly labeled diagram. Failure to meet this requirement will result in an automatic 1-point deduction. 


 Consider the equation 

 \[\left(3x^3 + xy^2 \right) \left(x^2y + 3y^3 \right) = (x-y)^7.\]


 	(a) Prove that there are infinitely many pairs $(x,y)$ of positive integers satisfying the equation. 


 	(b) Describe all pairs $(x,y)$ of positive integers satisfying the equation. 


 

	\item (\textit{USAJMO 2017}) Note: For any geometry problem whose statement begins with an asterisk ($*$), the first page of the solution must be a large, in-scale, clearly labeled diagram. Failure to meet this requirement will result in an automatic 1-point deduction. 


 ($*$) Let $ABC$ be an equilateral triangle and let $P$ be a point on its circumcircle. Let lines $PA$ and $BC$ intersect at $D$; let lines $PB$ and $CA$ intersect at $E$; and let lines $PC$ and $AB$ intersect at $F$. Prove that the area of triangle $DEF$ is twice the area of triangle $ABC$.


 

	\item (\textit{USAJMO 2017}) Note: For any geometry problem whose statement begins with an asterisk ($*$), the first page of the solution must be a large, in-scale, clearly labeled diagram. Failure to meet this requirement will result in an automatic 1-point deduction. 


 Are there any triples $(a,b,c)$ of positive integers such that $(a-2)(b-2)(c-2) + 12$ is a prime that properly divides the positive number $a^2 + b^2 + c^2 + abc - 2017$?


 

	\item (\textit{USAJMO 2017}) Note: For any geometry problem whose statement begins with an asterisk ($*$), the first page of the solution must be a large, in-scale, clearly labeled diagram. Failure to meet this requirement will result in an automatic 1-point deduction. 


 ($*$) Let $O$ and $H$ be the circumcenter and the orthocenter of an acute triangle $ABC$. Points $M$ and $D$ lie on side $BC$ such that $BM = CM$ and $\angle BAD = \angle CAD$. Ray $MO$ intersects the circumcircle of triangle $BHC$ in point $N$. Prove that $\angle ADO = \angle HAN$. 


 

	\item (\textit{USAJMO 2017}) Note: For any geometry problem whose statement begins with an asterisk ($*$), the first page of the solution must be a large, in-scale, clearly labeled diagram. Failure to meet this requirement will result in an automatic 1-point deduction. 


 Let $P_1, \ldots, P_{2n}$ be $2n$ distinct points on the unit circle $x^2 + y^2 = 1$ other than $(1,0)$. Each point is colored either red or blue, with exactly $n$ of them red and exactly $n$ of them blue. Let $R_1, \ldots, R_n$ be any ordering of the red points. Let $B_1$ be the nearest blue point to $R_1$ traveling counterclockwise around the circle starting from $R_1$. Then let $B_2$ be the nearest of the remaining blue points to $R_2$ traveling counterclockwise around the circle from $R_2$, and so on, until we have labeled all the blue points $B_1, \ldots, B_n$. Show that the number of counterclockwise arcs of the form $R_i \rightarrow B_i$ that contain the point $(1,0)$ is independent of the way we chose the ordering $R_1, \ldots, R_n$ of the red points. 


 

\end{enumerate}